
\section{Introduction}
An aircraft-inspection platform needs an airframe-relative pose to associate each observation with the inspected structure. The relevant decision combines geometry, uncertainty, available processing time and the provenance of the estimate. Completing a registration can improve accuracy, but the result may arrive after the observation request has expired. An intermediate estimate can arrive earlier, although a small residual alone does not establish that it is reliable. The engineering problem is to decide which result may be used, which computation remains worthwhile, and which statistical statement is applicable to that decision.

These decisions expose a gap between numerical evidence and its operational meaning. Two estimates from the same LiDAR scan share a prior and likelihood; they must not be fused as independent measurements. A risk bound calibrated for one operating mixture cannot be attached to a different aircraft, workload or event simply because both systems return a pose. A four-second localization event and the availability of a fresh observation pose within 0.3~s are different requirements. Representing these distinctions explicitly makes them executable rather than leaving them to interpretation of a scalar confidence score.

Engineering knowledge representations connect contextual facts, evidence and decisions \cite{wang2021process,hao2021knowledge,su2025manufacturing}. Contract-based system design already formalizes assumptions and guarantees, including timing and scheduling \cite{benveniste2018contracts}. Runtime-assurance architectures also separate a performance controller from the conditions under which its output is admissible \cite{mehmood2022simplex}. We build on these principles for a specific estimation interface: delayed, correlated versions of one registration result, coupled to distribution-specific statistical event cards. The contribution concerns the representation and use of this decision knowledge, together with its measurable effect on registration-linked observation requests.

RAVC-Loc retains one robust ICP computational state and one immutable source prior. We extend it with \emph{executable operating contracts}. Before registration, a declared task intersects the event-compatible policy sets and selects a frozen policy by calibration cost; an empty intersection produces explicit refusal. During registration, observed geometric quality, declared sensor settings and prior uncertainty govern admissible computational actions. A compact model forecasts several terminal iteration caps, allowing the scheduler to account for nonmonotonic solution quality. A predicted-work screen restricts refinement as the observation deadline approaches. Delivered candidates replace their source version and replay causal measured motion. Statistical event cards bind the selected policy and forecast, source semantics, input mixture, geometric artifact, installation model and requested event.

The study makes three contributions. First, it develops an executable representation that separates candidate eligibility, computational admissibility, source-version consistency and statistical authorization. Requested events change the compatible policy set and can cause policy substitution or refusal, rather than merely changing a confidence label. Second, it couples this representation to multi-action computation forecasts and a serial replacement/replay service; the analysis distinguishes estimator consistency, the value of early installation and the assumptions needed for finite-family risk control. Third, it tests these mechanisms with fresh independent episodes, matched CART and random-forest schedulers, a geometry excluded from forecast fitting, influence and tail analyses, and complete simulated routes. A secondary frozen-record analysis measures event-driven policy substitution and its observation/trajectory tradeoffs without counting reused records as new evidence.

The central empirical question is whether explicit contracts add useful decision capability beyond a numerical stopping score. The experiments accordingly distinguish the effect of rules from the effect of the forecast model and from the choice of operating price. They also report circumstances in which a cheaper fixed policy is preferable or the specified risk allowance cannot be certified. The original A19 and earlier A20 evidence remain available as separate, reproducible studies.

Figure~\ref{fig:contracts} summarizes how a declared inspection task, calibrated evidence and causal localization execution connect.
